% ===================================================================
%  4. TAULA — Heldutasuna eta Zahartzaroa.
%  Traduction 2026 d'apres texto/io, controlee sur texto/fr, texto/en
%  et texto/es. Consigne : voir texto/eu/10-taula-01.tex.
%
%  LE TABLEAU DE LA PARENTE, ET LE BASQUE Y DEMANDE PLUS DE SOIN QUE
%  TOUTE AUTRE COLONNE : la sœur ainee de Marthe est « ahizpa », parce
%  que Marthe est une femme ; le frere de la mariee est « neba » pour
%  la meme raison. Cette regle-la n'est pas un ornement de dialecte,
%  c'est la grammaire ordinaire.
%
%  LA RELATIVE BASQUE SE PREPOSE, et c'est ce qui menace le bloc
%  c2-07-1 : « medikuak idatzi berri duen errezeta » mettrait le
%  medecin (25) devant l'ordonnance (27), la ou l'ido donne 27 puis 25.
%  On la postpose donc en apposition — « errezeta (27), medikuak (25)
%  idatzi berri duena » — comme au tableau 3.
% ===================================================================

%%K t04-apar-1 apar
\VUsaut{4.0mm}
\VUtitre{15.0pt}{60}{4. TAULA}
\VUsaut{1.6mm}
\VUtitre{11.4pt}{0}{Heldutasuna eta Zahartzaroa.}
\VUsaut{3.2mm}

%%K t04-tit-1 sub
\VUcentre{11.4pt}{0}{\textit{Lehen agerraldia.}}
\VUsaut{1.6mm}
\VUtitre{11.4pt}{0}{Ezteia-bazkaria.}
\VUsaut{2.6mm}

%%K t04-tit-2 sub
\VUpk{10.2pt}{40}{Udazkena.}
\VUsaut{3.0mm}

%%K t04-c1-01-1 p
1. --- Gazteak hazi egin dira. Haietako bat, Joanes, ezkontzen ari da.
\VUgras{Ezteia-bazkarian} gaude, zeinak mahai berean bildu baititu senargaiaren
eta emaztegaiaren familiak.

%%K t04-c1-02-1 p
2. --- \VUgras{Udazkena} da, \VUgras{iraila}. \VUgras{Urriko} eta
\VUgras{azaroko} haize hotzak ez ditu oraindik landetako hosto berdeak eraman.
Arratseko airea epela eta usaintsua da; horregatik bazkaria \VUgras{berandaren}
\textsuperscript{(8)} azpian ematen da, lorategira begira.

%%K t04-c1-03-1 p
3. --- Urrunean, \VUgras{muinoaren} \textsuperscript{(3)} gainean, Joanesen
gurasoen etxea dago, \VUgras{jauregitxo} dotore bat \textsuperscript{(1)},
berdetasunean ezkutatua; ardo bikaina ematen duten mahasti ederrek estaltzen dute
muinoaren hegala. Mahastizainak zaindu eta maitatu egiten ditu.

%%K t04-c1-04-1 p
4. --- Mahatsondoen gainetik \VUgras{eper} talde bat \textsuperscript{(2)}
hegaldatzen da, \VUgras{basozainaren} \textsuperscript{(5)} hurbiltzeak izutua.
Eta hark, \VUgras{eskopeta} besapean eta \VUgras{ehiza-zorroa bizkarrean},
\VUgras{sasiak} \textsuperscript{(4)} miatzen ditu bere \VUgras{txakurrarekin}
\textsuperscript{(6)}. Jabetza zaintzen du eta ez die ehiza-lapurrei
\VUgras{ehiza} suntsitzen uzten.

%%K t04-c1-05-1 p
5. --- Kanpotik \VUgras{berandan} \VUgras{aihena} kiribiltzen da, zeinetik
\VUgras{mahats-mordo} astunak \textsuperscript{(7)} zintzilikatzen baitira.

%%K t04-c1-06-1 p
6. --- \VUgras{Mahaia} \textsuperscript{(16)} apaintzen duten udazkeneko lore
ederrak \VUgras{krisantemoak} \textsuperscript{(9)} dira; emaztegaiaren
\VUgras{aitak} \textsuperscript{(10)} bat sartu du frakaren begi-zuloan.

%%K t04-c1-07-1 p
7. --- Bazkaria amaitzera doa. \VUgras{Zerbitzaria} \textsuperscript{(11)}
postrearen platerak ekartzen hasi da eta laster jasoko ditu jada beharrezkoak ez
diren mahai-tresnak --- \VUgras{koilarak} \textsuperscript{(14)},
\VUgras{sardexkak} \textsuperscript{(12)}, \VUgras{aiztoak}
\textsuperscript{(13)}, \VUgras{azpilak} \textsuperscript{(15)}.

%%K t04-tit-3 sub
\VUpk{10.2pt}{40}{Familia.}
\VUsaut{3.0mm}

%%K t04-c1-08-1 p
8. --- Denek itxura alaia dute, eta senargaiaren \VUgras{amak}
\textsuperscript{(17)} bere seme-alabei begiratzen dien irribarreak erakusten du
zein zoriontsu dagoen.

%%K t04-c1-09-1 p
9. --- Ezkontza honek eskualdeko bi familia aberats batzen ditu. Joanes,
\VUgras{senarra} \textsuperscript{(18)}, lurjabe handi baten \VUgras{semea} da,
eta fabrikatzaile aurreratu baten \VUgras{suhi} bihurtzen da.

%%K t04-c1-10-1 p
10. --- \VUgras{Ezkontideak} gazteak dira, eta hala ere luzaroan
\VUgras{hitzemanak} egon dira. \VUgras{Emazte gaztea} \textsuperscript{(19)},
Marta, ederra dago bere \VUgras{soineko zurian}, ezkontza-lorez apaindua.

%%K t04-c1-11-1 p
11. --- Haren \VUgras{aitaginarreba} \textsuperscript{(20)} pozik dago
\VUgras{erran} hain on batekin, eta Joanesen \VUgras{amaginarrebak}
\textsuperscript{(21)} irribarre egiten du bere \VUgras{alaba} hain eder ikustean.

%%K t04-c1-12-1 p
12. --- Mahaiaren muturrean gainerako \VUgras{gonbidatuak}
\textsuperscript{(22)} daude eserita: \VUgras{ahaideak, lagunak, lehengusuak}.
\VUgras{Jantoki-nagusiak} \textsuperscript{(23)} serbileta besoan, bizkor
zerbitzatzen die.

%%K t04-tit-4 sub
\VUpk{10.2pt}{40}{Lagunak.}
\VUsaut{3.0mm}

%%K t04-c1-13-1 p
13. --- Mahaiaren eskuinaldean \VUgras{neskatxa} bat \textsuperscript{(24)} dago,
emaztegaiaren \VUgras{laguna}, eta gero emaztegaiaren neba, haren
\VUgras{ezkontza-mutila} \textsuperscript{(25)}. Bere \VUgras{ezkontza-neskatxarekin}
\textsuperscript{(26)} ari da hizketan, senargaiaren arreba, zeina ezkontza honen
bidez Martaren \VUgras{koinata} bihurtzen baita. Emakume gazte xarmangarria da.
\VUgras{Bitxi} ederrak ditu, \VUgras{perla-lepokoa} eta urrezko
\VUgras{eskumuturrekoa}.

%%K t04-c1-14-1 p
14. --- Haien ondoan zutitu eta kopa altxatu duen jauna familiaren
\VUgras{laguna} \textsuperscript{(27)} da. Senargaiaren \VUgras{lekukoetako} bat
da. \VUgras{Topa} egiten du, ezkonberrien \VUgras{osasunera} edaten du eta
zorion luzea opa die. Jantzi dotorez dago: fraka eta \VUgras{zapata lakatuak}
\textsuperscript{(28)}. Besotik darama \VUgras{amona} \textsuperscript{(29)},
zeina \VUgras{alarguna} baita.

%%K t04-c1-15-1 p
15. --- \VUgras{Frakez} \textsuperscript{(31)} jantzitako gizon gaztea, bibote
beltz finekoa, Joanesen \VUgras{koinatua} \textsuperscript{(30)} da. Ingeniari
ospetsua da. Oso \VUgras{dama} dotore baten \textsuperscript{(33)} ondoan dago
eserita, Martaren \VUgras{ahizpa nagusia} eta lore bat eskaintzera eta
\VUgras{arratsalde on} \VUgras{opatzera} lehengusuaren besotik datorren neskatxa
maitagarriaren ama. Dama honek \VUgras{belarritako} ederrak eta \VUgras{orrazkera}
dotorea ditu.

%%K t04-c1-16-1 p
16. --- Lehengusu txikia, bere \VUgras{blusan} \textsuperscript{(36)} eta
\VUgras{galtza zabaletan} \textsuperscript{(37)} hain barregarria, oso errukigarria
da. \VUgras{Umezurtza} \textsuperscript{(35)} da. Aita eta ama oso txikitan galdu
zituen. Haren osaba haren \VUgras{tutorea} \textsuperscript{(38)} da eta
ezkongabe geratu da haur gaixoa hazteko. Gizon ona da. Nahiko burusoila da eta
oso miopea; \VUgras{pintza-betaurrekoak} eramaten ditu.

%%K t04-tit-5 sub
\VUsaut{2.4mm}
\VUpk{10.2pt}{40}{Postrea.}
\VUsaut{3.0mm}

%%K t04-c1-17-1 p
17. --- Gonbidatuen atzean, \VUgras{mahai-zapiz} estalitako mahai batean,
\VUgras{postrea} \textsuperscript{(39)} dago jarrita. Ontzi handi batean fruta:
\VUgras{mertxikak} \textsuperscript{(40)}, \VUgras{udareak}
\textsuperscript{(41)}, \VUgras{sagarrak} \textsuperscript{(42)},
\VUgras{abrikotak} \textsuperscript{(43)} eta \VUgras{mahatsa}
\textsuperscript{(44)}. Ondoan \VUgras{ananak} \textsuperscript{(45)},
\VUgras{tarta} eder bat \textsuperscript{(46)}, \VUgras{likore-botilak}
\textsuperscript{(48)} eta \VUgras{xanpainarenak} \textsuperscript{(49)}, eta
\VUgras{plater} garbien piloak \textsuperscript{(47)}.

%%K t04-c1-18-1 p
18. --- \VUgras{Ardandariak} \textsuperscript{(50)} ardo zaharreko
\VUgras{botila} bat \textsuperscript{(52)} irekitzen du \VUgras{kortxo-kentzekoaz}
\textsuperscript{(51)}. \VUgras{Kortxoa} \textsuperscript{(53)} oso estu doa,
dirudienez. Ardandari honek \VUgras{serbileta} \textsuperscript{(54)} du besoan.

%%K t04-c1-19-1 p
19. --- Bazkariaren ondoren gizonak erretzeko gelara joango dira, eta damak
dantzaldirako janztera.

%%K t04-tit-6 sub
\VUsaut{5.0mm}
\VUcentre{11.4pt}{0}{\textit{Bigarren agerraldia.}}
\VUsaut{1.6mm}
\VUpk{11.4pt}{40}{Aitonaren izen-eguna.}
\VUsaut{2.6mm}

%%K t04-tit-7 sub
\VUpk{10.2pt}{40}{Bisita.}
\VUsaut{3.0mm}

%%K t04-c2-01-1 p
1. --- Hamar urte igaro dira. Joanesen gurasoak zahartu egin dira eta samurki
maite dituzte beren hiru bilobak: bi \VUgras{biloba} \textsuperscript{(2)} eta
\VUgras{biloba neska} bat \textsuperscript{(3)}. Gaur \VUgras{aitonaren
izen-eguna} denez, haurrak zoriontzera etorri dira.

%%K t04-c2-02-1 p
2. --- Petri txikia oso txikia da oraindik; hiru urte baino ez ditu.
\VUgras{Katua} \textsuperscript{(1)} beregana datorrela ikusten du. Haren
\VUgras{ile} zetatsu ederra eta \VUgras{buztan} luzea laztandu nahi ditu; ez zaio
burutik pasatzen \VUgras{hanka} dotore horien azpian atzapar zorrotz gaiztoak
daudela, eta harramazka egin diezaioketela.

%%K t04-c2-03-1 p
3. --- Haren ahizpa Gabriela, zeinak eskutik heltzen baitio, askoz serioagoa da,
eta Martzel, bere \VUgras{beroki} handiarekin \textsuperscript{(4)} eta larruzko
\VUgras{gerrikoarekin} \textsuperscript{(5)}, erabat helduaren itxura du, galtza
motzen azpitik bere \VUgras{galtzerdi luzeak} \textsuperscript{(6)} ageri badira
ere.

%%K t04-c2-04-1 p
4. --- Haien aitak neguko \VUgras{beroki} lodia \textsuperscript{(7)} jantzi du,
\VUgras{larruzko lepoarekin} \textsuperscript{(8)}, eta \VUgras{poltsikoan}
\textsuperscript{(9)} bufanda darama. Haren emazteak feltrozko kapela darama
\VUgras{lumekin} \textsuperscript{(10)}, \VUgras{jaka} \textsuperscript{(11)},
larruzko \VUgras{boa} \textsuperscript{(12)} eta \VUgras{eskuberogailua}
\textsuperscript{(13)}.

%%K t04-c2-05-1 p
5. --- Oso hotz egiten du, negua etorri baita. \VUgras{Elurra}
\textsuperscript{(19)} egun osoan aritu da, eta etxeen teilatuak erabat zuri
daude. \VUgras{Gauarekin} izotza gogortu egiten da are gehiago. Zeruan
\VUgras{ilargia} \textsuperscript{(17)} eta \VUgras{izarrak}
\textsuperscript{(18)} distiratzen dute, \VUgras{iluntasuna} zulatuz.

%%K t04-tit-8 sub
\VUsaut{4.2mm}
\VUpk{10.2pt}{40}{Gaixotasuna.}
\VUsaut{3.0mm}

%%K t04-c2-06-1 p
6. --- \VUgras{Itsu} gaixoa \textsuperscript{(20)}, zeina \VUgras{botikaren}
\textsuperscript{(16)} aurreko plaza zeharkatzen baitu bere \VUgras{pudelak}
\textsuperscript{(21)} gidatuta, oso zorigaiztokoa da. Justina,
\VUgras{neskamea} \textsuperscript{(14)}, sukaldetik \VUgras{katilu} bat
\textsuperscript{(15)} belar-ur oso bero dakar, eskuzabal gozotua.

%%K t04-c2-07-1 p
7. --- \VUgras{Amona} \textsuperscript{(23)}, zeinak mahaitik bere
\VUgras{betaurrekoak} \textsuperscript{(26)} hartu nahi baititu
\VUgras{errezeta} \textsuperscript{(27)} irakurtzeko --- \VUgras{medikuak}
\textsuperscript{(25)} idatzi berri duena ---, besaulkitik altxatzen da, bere
\VUgras{makilan} \textsuperscript{(24)} bermatuta.

%%K t04-c2-08-1 p
8. --- Sarritan \VUgras{gaitz} txikiak jasaten ditu, \VUgras{zorabioak},
\VUgras{hotzeriak} eta \VUgras{hortzetako mina}; hankak ahul ditu eta begiak ere
ez oso onak. Eta hala ere pozik dago bere \VUgras{mediku} zaharrari adarra
jotzeaz, zeinak beti agindu nahi baitio \VUgras{nahaste-edari} bat
\textsuperscript{(28)}. Gaur oso zoriontsu dago senide gazteak inguruan izateaz.

%%K t04-c2-09-1 p
9. --- Ondo itxitako gelan gozo dago. Marmolezko \VUgras{tximinian}
\textsuperscript{(29)} \VUgras{su ederra} \textsuperscript{(30)} dabil, eta gozo
dago bihotza piztu berri den \VUgras{lanpararen} \textsuperscript{(31)}
\VUgras{argi} epelean.

%%K t04-tit-9 sub
\VUsaut{4.2mm}
\VUpk{10.2pt}{40}{Aitona.}
\VUsaut{3.0mm}

%%K t04-c2-10-1 p
10. --- \VUgras{Aitonak} \textsuperscript{(33)} berak ere ahaztu du gaixorik egon
dela, \VUgras{erreumak} \VUgras{besaulkira} \textsuperscript{(34)} lotuta
daukala eta atzo bertan \VUgras{sukarra eta konorte-galtzeak} izan zituela.
Jada ez du gogoan iazko neguan \VUgras{birikeria} izan zuela eta \VUgras{eztul}
gogor eta mingarri batek asko sufriarazi ziola.

%%K t04-c2-10-2 p
Eta hala ere nekez sendatu zen. Orain hobeto dago. Baina oraindik min ematen dio
\VUgras{burkoaren} \textsuperscript{(38)} gainean daukan hankak; eta burkoa
\VUgras{aulkitxo} baxu baten \textsuperscript{(39)} gainean jarrita dago.

%%K t04-c2-11-1 p
11. --- Bere \VUgras{txabusina} epelean \textsuperscript{(35)} kontuz bilduta
dagoenean, nakarrezko \VUgras{botoi} handiekin \textsuperscript{(36)}, eta haren
ondoan, egunkaria irakurtzeko, \VUgras{umezurtz} txikia \textsuperscript{(40)}
eserita dagoenean, \VUgras{txano} zuriarekin, \VUgras{soineko} urdinarekin
\textsuperscript{(42)} eta \VUgras{kaparrarekin} \textsuperscript{(41)}, ez da
kexatzen.

%%K t04-c2-12-1 p
12. --- Urtero, \VUgras{egutegiak} \textsuperscript{(43)} \VUgras{abenduaren}
lehen \VUgras{eguna} berriro dakarrenean, irrikaz itxaroten die bere
\VUgras{bilobei}, badaki eta ez dutela \VUgras{egun} garrantzitsu hori ahaztuko.

%%K t04-c2-13-1 p
13. --- Hunkituta gogoratzen ditu urruneko egunak, berak enperadorearen
\VUgras{mariskal} ospetsua zoriontzera joaten zenekoak, zeinaren
\VUgras{erretratua} \textsuperscript{(44)} horman baitago zintzilik eta zeina
haurren \VUgras{birraitona} baita.
\VUsaut{6.0mm}
\VUfilet{22mm}
